% Handwritten interpretation; numerical values are generated separately.
\subsection{Compression preserves the ordinary dense cut at modest rank}

The first question is whether a compact factor represents the cut values that
the solver is intended to improve. Across six full development datasets, three
factor seeds and five partition perturbations per factor, the largest rank-256
error for the ordinary kernel is \PlainFidelityError{} of its off-diagonal
similarity mass. Error decreases with rank (Figure~\ref{fig:powered}). Thus, in
this collection, a modest factor supports accurate evaluation of the dense cut
without storing its full matrix. This tests objective fidelity across the
examined partitions; optimization quality remains a separate question.

For comparison, the top-32 graph retains a mean \GraphMassDefault\% of the
ordinary kernel's similarity mass on the same six datasets. Raising its degree
to 128 increases this to \GraphMassHigh\%. Compression and sparsification thus
retain different aspects of the dense objective. The graph-degree study and
its returned-partition scores appear in the supplement.

\begin{figure}[htbp]
\centering\includegraphics[width=6.5in]{revision/powered_fidelity.pdf}
\caption{Cut-value fidelity for the ordinary and cubed kernels. Each point is
the largest absolute factor-versus-exact cut difference, divided by the exact
kernel's off-diagonal similarity mass, across six full datasets, three seeds
and five partition perturbations. The powers have different denominators.
The cubed kernel requires substantially more rank.}
\label{fig:powered}
\end{figure}

\subsection{Local similarity measures a distinct aspect of separation}

A faithful dense objective is useful only if its scientific meaning is clear.
On synthetic datasets with known families, local leakage correlates with
cross-fold family contamination at Spearman $\rho=\SyntheticLocalRho$, compared
with $\SyntheticGlobalRho$ for the global ratio on the same \SyntheticCount{}
observations. Pooled correlations span several constructed regimes; their
individual results appear in the supplement. The continuum regime favors the
global score, whereas four other regimes favor the local score. This supports measuring
close-neighbour contamination separately from total overlap.

Predictions provide a second interpretation. In the random and plain-low-rank
D-MPNN partitions, similarity to the nearest training molecule correlates
negatively with absolute ensemble error on ESOL, FreeSolv and Lipophilicity
($\rho=\LocalEsolRho$, $\LocalFreesolvRho$ and $\LocalLipoRho$, respectively).
On Lipophilicity, the per-point global score has a weak positive correlation
($\GlobalLipoRho$). BBBP is an exception to the local trend
($\rho=\LocalBbbpRho$), and within-split associations vary. Local similarity
therefore captures an empirically relevant aspect of prediction, with a
model- and dataset-dependent relationship to error. Within-split analyses and
prediction provenance are supplied in the supplement.

\subsection{Power and grouping control local overlap, with a fidelity cost}

The second question is whether the same framework can reduce close-analogue
contamination. We varied power over $\{2,3,4\}$ and grouping threshold over
$\{0.3,0.4,0.5,0.6\}$ on six full datasets, using exact fold counts and three
seeds. Plain, power-only and group-only controls complete the 306-partition
study. The resulting global--local frontier is supplied in the supplement.
At the fixed combined setting, local leakage is lower than plain low rank on
all six datasets under each of three Morgan ramps, the top-five-neighbour ramp
and the fraction of test molecules above similarity 0.6. The improvement
therefore extends beyond the ramp whose floor coincides with the grouping
threshold. MACCS, used only for evaluation, shows lower mean nearest-training
similarity on \RobustMaccsWins{} of six datasets. Its dataset-bootstrap interval
for the mean paired change includes zero. Local control is consequently strongest
in the representation used to construct the split.

That control has a computational trade-off. For the cubed kernel, the maximum
rank-256 cut error is \PowerFidelityError{} of powered mass, falling to
\PowerHighRankError{} at rank 512 (Figure~\ref{fig:powered}). The ordinary
kernel's fidelity therefore cannot be assumed after sharpening. The normalized
spectral estimates are often above one and give a vacuous uniform bound.
Nevertheless, combined-setting local leakage changes less with rank: BACE's
mean local ratio falls from \BaceRankSmall{} at rank 32 to
\BaceRankDefault{} at rank 256 and \BaceRankLarge{} at rank 512.
Grouping can preserve local structure even when the powered factor is coarse.
We therefore evaluate the returned partition as well as the factor's objective
error. Full rank curves and alternative metrics appear in the supplement;
accompanying result records retain group-relaxation counts.

\subsection{Full-size baseline comparison}

\IfFileExists{benchmark_numbers.tex}{The full benchmark separates the benefit of compressing a dense objective from
the effect of changing that objective. At tolerance 0.01, plain low rank has
lower global overlap than the quality-preset graph comparator on
\BenchPlainGraphGlobalWins{} of \BenchPlainGraphGlobalN{} molecular datasets
(Figure~\ref{fig:benchmark-global}). The combined setting reduces local leakage
relative to plain low rank on \BenchCombinedPlainLocalWins{} of
\BenchCombinedPlainLocalN{} datasets, while raising the ordinary dense cut.
Its local score is lower than the graph comparator on
\BenchCombinedGraphLocalWins{} of \BenchCombinedGraphLocalN{} datasets
(Figure~\ref{fig:benchmark-local}). The mean combined-minus-graph local
change is \BenchCombinedGraphLocalChange{}, reflecting the differing
dataset effects. These comparisons expose the choice the
framework offers: the plain factor directly targets distributed overlap;
power and grouping spend some of that global advantage on local separation.
The accompanying machine-readable tables include both tolerances, all reference methods, realized
counts and group relaxations. Paired dataset effects, bootstrap intervals and
Holm-adjusted sign tests are supplied in the supplement.

\begin{figure}[!htbp]
\centering\includegraphics[width=6.5in]{revision/benchmark_global.pdf}
\caption{Global overlap on eleven full molecular datasets at size tolerance
0.01, ordered by dataset size. Points are means of three split seeds; bars show
their standard deviations. The dotted line is the count-matched random
expectation. All methods retain the same input entities. The graph comparator
uses exact top-32 construction, an ordinary pairwise Graph API, the quality
preset and CUT objective, with lower/upper count repair. The supplement reports
all methods, actual counts, degree sensitivity and the two materials checks.}
\label{fig:benchmark-global}
\end{figure}

\begin{figure}[!htbp]
\centering\includegraphics[width=6.5in]{revision/benchmark_local.pdf}
\caption{Test-only nearest-training-neighbour leakage for the same partitions
and scoring protocol as Figure~\ref{fig:benchmark-global}. Ratios are normalized
by count-preserving random permutations. Bars show split-seed standard
deviations, truncated at zero where necessary. Controlling total overlap and
close-neighbour contamination gives different method orderings.}
\label{fig:benchmark-local}
\end{figure}

DataSAIL provides a closer comparison for the dense objective. Its saved
assignments cover seven datasets, with source outcomes classified separately
from missing attempts and resource-policy exclusions. At their exact native
counts, plain low rank has lower global overlap on
\BenchPlainDataSailGlobalWins{} of \BenchPlainDataSailGlobalN{} datasets.
The combined setting has lower global overlap on
\BenchCombinedDataSailGlobalWins{} of \BenchCombinedDataSailGlobalN{} and lower
local leakage on \BenchCombinedDataSailLocalWins{}, with
\BenchCombinedDataSailLocalTies{} tie. Thus compression provides a practical
dense-cut objective, while the local setting remains a choice among competing
criteria. Reused assignments are rescored with the common protocol and supply
no new runtime.

The published local and cluster methods clarify the retention contract.
Lo-Hi satisfies its 0.4 cross-fold threshold on all seven eligible datasets,
retaining between 98.0\% and 99.9\% of their molecules. Low-rank controls match
each retained subset and its exact counts; their capped groups offer soft local
control rather than that hard threshold guarantee. In the transferred UMAP
recipe, seven native cluster holdouts and their count-matched controls retain
every molecule. The supplement reports both studies separately. The FreeSolv native UMAP
failure and MUV memory-policy exclusion remain unscored.

\IfFileExists{runtime_numbers.tex}{
\subsection{Dense splitting has a distinct computational advantage}

The same-host comparison measures the public split call at exactly the native
DataSAIL counts. Across four cases, median combined-setting times range from
\RuntimeCombinedMinSeconds{} to \RuntimeCombinedMaxSeconds{} seconds;
DataSAIL takes \RuntimeDataSailMinSeconds{} to
\RuntimeDataSailMaxSeconds{} seconds, including distance construction,
clustering and its integer program (Figure~\ref{fig:runtime-matched}).
Local quality remains dataset-dependent: on BACE the combined mean local ratio
is \RuntimeBaceCombinedLocal{}, compared with
\RuntimeBaceDatasailLocal{} for DataSAIL. Speed and the desired separation
criterion must therefore be considered together. Solver outcomes, initialization
policy, individual timings and all exact counts appear in the supplement.

\begin{figure}[!htbp]
\centering\includegraphics[width=6.5in]{revision/runtime_matched.pdf}
\caption{CPU cost and local overlap on four same-host cases at exactly matched
DataSAIL fold counts. Colour denotes method and shape denotes dataset.
Low-rank points average three split seeds for local leakage and use median
split-call time; DataSAIL has one solve per case. The clock excludes feature
loading and scoring. DataSAIL's distance construction, clustering and solver
are included; low-rank calls exclude a 100-entity warm-up. Solver-reported
optima and time-limited feasible outcomes are distinguished in the supplement.}
\label{fig:runtime-matched}
\end{figure}
}{}
}{
\noindent\textbf{Revision in progress:} the full-size, version-stamped baseline
rerun is still running. The final comparison will use the common protocol in
Methods. No earlier runtime is used here to claim a speedup, and no partial grid
is presented as a completed thirteen-dataset result.}

\subsection{Computational scaling differs between settings}

On an H100 NVL, plain low rank partitions two million 2048-bit ChEMBL
fingerprints in \ScalePlainLargeSeconds{} seconds, with
\ScalePlainLargeMemory{} GiB of peak GPU allocation. The measurement includes
the public split call's factor construction, four restarts and refinement,
after a small warm-up; loading features and evaluating leakage are outside
the clock. The combined setting takes \ScaleCombinedLargeSeconds{} seconds
at 100,000 entities using exact tiled neighbours. These settings have different
computational profiles (Figure~\ref{fig:scaling-revision}). The combined
pipeline is measured through 100,000 entities, and its performance is not
extrapolated from the two-million-entity plain-factor result. Per-case memory,
settings and output counts are supplied in the supplement.

\begin{figure}[htbp]
\centering\includegraphics[width=6.5in]{revision/scaling.pdf}
\caption{Split-call time on nested ChEMBL samples using the standard solver,
rank 256 and size tolerance 0.01. All cases use H100 NVL hardware with two pinned
CPU cores per worker. Each point is a fresh process with a 100-entity warm-up;
feature loading and quadratic diagnostics are excluded. The plain and combined
settings have separate measured ranges. k denotes thousands and M millions.
Full per-case timings, GPU peaks and host RSS appear in the supplement.}
\label{fig:scaling-revision}
\end{figure}

\subsection{Small-assay separation is strongly sensitive to target composition}

On all ten prespecified external assays, the combined setting yields mean
test ROC-AUC \ExtCombinedRF{} for RF and \ExtCombinedXGB{} for boosted trees.
Random splitting gives \ExtRandomRF{} and \ExtRandomXGB{}, respectively;
the published SIMPD references give \ExtSimpdRF{} and \ExtSimpdXGB{}.
In this assay collection SIMPD's aggregate scores are close to random, so a
large drop from SIMPD is not by itself evidence of realistic project transfer.
Several combined partitions produce below-chance AUC. Absolute scores and fold class proportions are supplied in the supplement.

At unchanged fold counts, matching class proportions raises the combined
setting's mean AUC to \ExtCombinedBalancedRF{} for RF and
\ExtCombinedBalancedXGB{} for boosted trees (Figure~\ref{fig:external-revision}).
The intervention changes a median \ExternalMoved\% of the combined assignments.
Mean standardized target shift falls from \ExternalShiftOriginal{} to
\ExternalShiftBalanced{}. Plain low rank is also sensitive, changing from
\ExtPlainRF{} to \ExtPlainBalancedRF{} for RF. The intervention simultaneously
changes which chemical series appear in each fold, so it is not a causal
estimate of the contribution of class proportions alone. Nevertheless, the
original chance-level results cannot be interpreted simply as successful
construction of a more useful benchmark.

\begin{figure}[htbp]
\centering\includegraphics[width=6.5in]{revision/external_rf.pdf}
\caption{Absolute RF test AUC on ten external assays before and after the
label-balance sensitivity intervention. Each assay is weighted equally after
averaging split and model seeds. Error bars are descriptive 95\% bootstrap
intervals over assay means. The intervention uses target labels, preserves
fold counts and can break structural groups. It is not part of the label-blind
method. SIMPD supplies one published split per assay. The dotted line is AUC
0.5. Boosted-tree results and class proportions appear in the supplement.}
\label{fig:external-revision}
\end{figure}

These assays delimit the purpose of stronger structural separation. A
new-series test can also separate the activity distributions associated with
those series, as anticipated by the distinction between the Lo and Hi tasks
\citep{steshin2023lohi}. The combined setting is useful when that transfer is the
intended question and sufficient predictive signal remains. For within-series
interpolation, its imposed boundary can be excessive. Structural leakage,
response composition and absolute model scores are needed together to assess
the resulting benchmark.
